% Programming with InitechBase, Chapter 1: Programming in a Nutshell
% (AN-16; annex Sec 8.2 P1).  Panels replayed: examples/p1-*.dot (dot-prompt
% route) and ../../examples/samir/salary (the program route of
% harness/pubs/samir_replay.c, samir_do).  Sources samir_main.c
% repl_try_do_file, cmd/proc.c, cmd/flow.c.
\chapter{Programming in a Nutshell}
\Gates{test-pubs-examples, test-interp-proc, test-interp-flow}
A program is a list of commands written in a file.  When you run the program,
InitechBase carries out the commands one after another, as if you were typing
them at the dot prompt, but faster and without mistakes.  In this chapter you
see what a program is made of, run the first program of the Initech Accounting
System, and read it line by line.

\begin{ChapterCovers}
\item From the dot prompt to a program
\item Program files
\item Running a program with \cmd{DO}
\item Reading the first program of the system
\end{ChapterCovers}

\ATEHead{Preparing for This Chapter}
Copy the Sample Programs and Utilities diskette to the disk you work on, as
\doctitle{Getting Started} describes, and start InitechBase from it.  Copy the
diskette again before each chapter: the programs of later chapters change the
ledger.

\ATEHead{From the Dot Prompt to a Program}
At the dot prompt you can look at the ledger an invoice at a time:
\begin{Steps}
\item Type \type{USE INVOICE}, then \type{GO TOP}.
\item Type \type{? INVNO, CUST, AMOUNT}, then \type{SKIP}, then \type{? INVNO, CUST,
AMOUNT}.
\end{Steps}
\DotExample{examples/p1-dot}
To look at every open invoice this way you would type the same two commands
for each record, and decide each time whether the invoice is paid.  A program
does the typing and the deciding for you.

\ATEHead{Program Files}
A program file is a text file with the extension \file{.PRG}, with one
InitechBase command on each line.  You write it with any editor that saves
plain text, such as the InitechDOS line editor \file{EDLIN}.  The commands are
those you use at the dot prompt, with a few more that exist only in programs:
\cmd{IF}, \cmd{DO CASE} and \cmd{DO WHILE}, which decide and repeat, and
\cmd{PROCEDURE} and \cmd{RETURN}, which divide a program into parts.

A program uses the databases that are open when it runs.  Open them at the dot
prompt with \cmd{USE} before you run the program: \cmd{USE}, \cmd{CLOSE} and
\cmd{QUIT} are dot prompt commands, and a program does not use them.

Here is a very small program.  It sets the decimal places, stores a salary in
a memory variable and displays it with an allowance added:
\VerbatimInput[fontsize=\small,xleftmargin=4mm]{../../examples/samir/salary.prg}
When it runs, it displays:
\begin{Transcript}{../../examples/samir/salary}
\end{Transcript}

\ATEHead{Running a Program}
\cmd{DO} runs a program.  You need not type the extension.  The first program of
the Initech Accounting System, \file{OPENINV.PRG}, lists the invoices that have
not been paid:
\begin{Steps}[resume]
\item Type \type{USE INVOICE}.
\item Type \type{DO OPENINV}.
\ATEResponse{The program lists the open invoices and counts them.}
\end{Steps}
\DotExample{examples/p1-openinv}
When the program ends, InitechBase shows the dot prompt again.

\ATEHead{Reading the Program}
\ProgramListing{OPENINV.PRG}
\SideHead{Comments}The first lines begin with an asterisk.  They are comments:
InitechBase passes over them.  Every program of the system begins with its
name, its purpose and the database it expects.

\SideHead{Headings}Two \cmd{?} commands print the title and the column
headings.

\SideHead{The count}\texttt{MCOUNT = 0} creates the memory variable that counts
the open invoices.

\SideHead{The loop}\cmd{GO TOP} moves to the first record.  The commands
between \cmd{DO WHILE} and \cmd{ENDDO} are repeated as long as the end of the
file has not been reached; \cmd{SKIP}, at the end of the loop, moves to the
next record each time.

\SideHead{The decision}Inside the loop, \cmd{IF .NOT. PAID} chooses the open
invoices.  For each, \cmd{?} prints the invoice number, the customer and the
amount, and the count grows by one.  \fn{STR()} turns the amount into
characters, so that it can be joined to the other values with \texttt{+}.

\SideHead{The end}After the loop, the program prints the count and
\cmd{RETURN} ends it.

\ATEHead{Where Do You Go From Here?}
Chapter~2 looks more closely at the structures that decide and repeat.

\ATEHead{Summary}
\begin{enumerate}
\item A program is a text file of commands, with the extension \file{.PRG}.
\item Open the databases with \cmd{USE}, then run the program with \cmd{DO}.
\item Comment lines begin with an asterisk.
\item \cmd{DO WHILE .NOT. EOF()} \ldots{} \cmd{SKIP} \ldots{} \cmd{ENDDO} visits
  every record; \cmd{IF} chooses records.
\item \cmd{RETURN} ends the program.
\end{enumerate}
